%%=============================================================================
%% Samenvatting
%%=============================================================================

% TODO: De "abstract" of samenvatting is een kernachtige (~ 1 blz. voor een
% thesis) synthese van het document.
%
% Een goede abstract biedt een kernachtig antwoord op volgende vragen:
%
% 1. Waarover gaat de bachelorproef?
% 2. Waarom heb je er over geschreven?
% 3. Hoe heb je het onderzoek uitgevoerd?
% 4. Wat waren de resultaten? Wat blijkt uit je onderzoek?
% 5. Wat betekenen je resultaten? Wat is de relevantie voor het werkveld?
%
% Daarom bestaat een abstract uit volgende componenten:
%
% - inleiding + kaderen thema
% - probleemstelling
% - (centrale) onderzoeksvraag
% - onderzoeksdoelstelling
% - methodologie
% - resultaten (beperk tot de belangrijkste, relevant voor de onderzoeksvraag)
% - conclusies, aanbevelingen, beperkingen
%
% LET OP! Een samenvatting is GEEN voorwoord!

%%---------- Nederlandse samenvatting -----------------------------------------
%
% TODO: Als je je bachelorproef in het Engels schrijft, moet je eerst een
% Nederlandse samenvatting invoegen. Haal daarvoor onderstaande code uit
% commentaar.
% Wie zijn bachelorproef in het Nederlands schrijft, kan dit negeren, de inhoud
% wordt niet in het document ingevoegd. 

\IfLanguageName{english}{%
\selectlanguage{dutch}
\chapter*{Samenvatting}
\selectlanguage{english}
}{}

wat, waarom, hoe onderzoek gedaan, resultaten, conclusie en relevantie voor het werk?

Deze bachelor thesis categoriseert en interpreteerd buffer pool System Management Facility (SMF) records van het product IBM Message Queue op een reproduceerbare, betrouwbare 
en deterministische manier. Om een globaal beeld te geven van het gezondheid van buffer pools.

Buffer pools zijn in geheugen caches die helpen met het normaliseren en balanceren van kostbare IO operaties door veel gebruikte data bij te houden. Het configureren
van buffer pools is belangrijk om een veilige, beschikbare, bestendige IBM MQ infrastructuur te creëren. Buffer pools creeren SMF records die verschillende data 
bevat over het gebruik en statistieken van dat buffer pool. Het is dus een rijke bron van informatie. 


SMF records zijn echte wel complexe data records in zowel hoeveelheid, types, subtypes. Daarnaast heeft de huidige it landschap, 
vakkennis, andere factoren in elk bedrijf invloed interpretatie van SMF. De resulterende interpretaties kunnen verschillen 
door ongelijke kennis, afspraken, achtergrond. 

De hoofdonderzoeksvraag dat word beantwoord is of het mogelijk is om SMF records te categoriseren en of het leid tot kennis en interpretatie waarmee
gehanteerd kan worden vergeleken met manuele interpretatie en validatie. 

Het onderzoek bestaat uit het creëren van een categorisatie framework dat gevisualizeerd word door een PoC. Het categorisatie framework is 
ontworpen op een regel gebaseerde ontwerp om zich te houden aan deterministische, betrouwbare en verklarende redenen. 

Daarnaast is er een PoC gecreeërd om het classificatie framework op een overzichtelijke manier te visualizeren. 

Het onderzoek resulteert in een verklaarbare, deterministische, betrouwbare manier om SMF records te categorizeren en interpreteren. Het is echter niet 
perfect, doordat buffer pools een kleine deel zijn van een groter netwerk van IBM MQ subsystemen. Het PoC kan verfijnd en uitgebreid worden en 
verbonden worden met SMF records van andere IBM MQ subsystemen zoals: Channels, Queue Managers en Coupling facility om relaties te vinden.    

Deze bachelor thesis concludeert dat het mogelijk is om bepaalde inzichten op een betrouwbare, deterministische en verklaarbare te verkrijgen. De inzichten zijn een klein deel van het 
gehele situatie en is niet volledig representatief. Het is echter een goeie startpunt om richting te geven in mogelijke onderzoek naar de bron van slechte performantie. 



%%---------- Samenvatting -----------------------------------------------------
% De samenvatting in de hoofdtaal van het document

\chapter*{\IfLanguageName{dutch}{Samenvatting}{Abstract}}


